Studying how the trajectories of samples evolve over time in a setting where we have access to snapshots of marginals over time is important in many fields, including in single-cell biology or climate science. The central difficulty is that observations are cross-sectional: samples are often unpaired across marginals, intermediate trajectories are unobserved, and high dimensional measurements often lie along a curved low-dimensional manifold, where standard Euclidean distances can be unreliable. This gives rise to two key considerations and active areas of research: handling the discrete marginal structure of the data and explicitly using the Riemannian manifold-structure of the dataset as a whole.

The main frameworks that have appeared for solving this trajectory inference problem structure are Optimal Transport (OT) and Flow Matching (FM) \cite{lipmanFlowMatchingGenerative2023, liuRectifiedFlowFlow2022}. Flow matching is one common training framework used to learn a vector field that defines an ordinary differential equation (ODE) which can be integrated using an ODE solver to predict flows from source to target distributions. In the trajectory inference problem, it is not just matching the source distribution that is valuable, but also predicting interpolations between marginal distributions that are logical and usable for downstream applications, like causality inference \cite{sunCFlowsRevealingDynamic}. Optimal Transport is used to pair samples in adjacent marginals to stabilize training and increase inference speed but commonly relies on pairing samples to minimize total Euclidean distance (2-Wasserstein) which ignores the manifold-structure of high-dimensional data \cite{tongImprovingGeneralizingFlowbased2024}.

Separately, methods have been develop to interrogate the connectivity and geometric structure of graphs, using the graph Laplacian \cite{belkinLaplacianEigenmapsDimensionality2003}. When the graph used is a k-Nearest Neighbors graph of data in high dimensions, this can produce a graph and laplacians that encode spectral information about the lower-dimensional manifold and geometric relationship (distances) between samples \cite{lipmanBiharmonicDistance2010, moonPHATEVisualizingStructure2019}.

\paragraph{Spectral-Cost Multi-Marginal Flow Matching (\ourmodel)} The core motivation for our development is to inject information about the manifold-structure of the dataset into Flow Matching while maintaining simulation-free training, which is difficult in existing manifold-based FM models, such as Riemannian Flow Matching \cite{chenFlowMatchingGeneral2024a}. One knob that we can change while maintaining a simulation-free FM training is the optimal transport ground costs used to pair samples from adjacent marginals for training. Ambient euclidean distances can be noisy or unreliable on high-dimensional datasets, motivating the use of a cost function for optimal transport pairings that will respect the distances between data points on the ambient data manifold. A graph Laplacian spectral distance can encode manifold connectivity and branching structure, encouraging OT couplings to respect the geometry of the observed data. We introduce \ourmodel, which uses kNN-graph-Laplacian-based spectral costs in OT pairings in the multi-marginal Flow Matching setting.

\paragraph{Related Work} Two models that learn a vector field perform trajectory inference, but are based on Neural ODEs and not Flow Matching are TrajectoryNet \cite{tongTrajectoryNetDynamicOptimal2020} and MIOFlow \cite{huguetMIOFlowManifoldInterpolating2022a}. More recent versions that use Neural ODEs are GAGA \cite{sunGAGAGeometryAwareGenerative2025} and CFlows \cite{sunCFlowsRevealingDynamic} (all out of Prof. Krishnaswamy's lab at Yale). On the Flow Matching side, MMFM \cite{rohbeckMODELINGCOMPLEXSYSTEM2025} is one of the earlier examples of multi-marginal flow matching. Separately, work has been done to use Riemannian manifolds in Flow Matching, like in Metric Flow Matching \cite{kapusniakMetricFlowMatchinga} and Riemannian Flow Matching (RFM) \cite{chenFlowMatchingGeneral2024}. \ourmodel sits somewhere between MMFM and RFM, using information about the manifold-structure of the data to augment multi-marginal Flow Matching.

\paragraph{Relationship to current work} This project is related to another manuscript, which has been accepted to ICML 2026, but is not yet publicly available \cite{kansalOTPFM2026}, in which I am the second author. That work develops multi-marginal flow matching with optimal-transport potentials whereas we look at computing graph spectral distance costs for OT couplings in this work.